% results
\subsection*{TreeData stores lineage trees alongside other single-cell measurements}

To integrate lineage trees into standard single-cell workflows, we developed TreeData, an extension of AnnData that adds support for storing trees (Figure~2A). TreeData adds two fields, \code{obst} and \code{vart}, which store trees as \code{networkx} directed graphs \citep{hagberg2008networkx} aligned with observations (cells) and variables (genes), respectively. Node and edge attributes, such as node time, are stored as key-value pairs on the tree. The AnnData interface (\code{X}, \code{obs}, \code{var}, \code{obsm}, \code{varm}, \code{layers}, \code{uns}, etc.) remains unchanged, and because TreeData is a subclass of AnnData, it can be passed directly to Scanpy, Squidpy, and other scverse tools. TreeData objects can be saved to disk as \code{.h5td} or \code{.zarr} files, with I/O times comparable to AnnData, particularly for scRNA-seq data where expression matrices dominate I/O (Figure~S1A).

TreeData supports three \code{alignment} options (Figure~2B): (1) \code{"leaves"} indicates that observations correspond to the leaves of the tree, as is typically the case for molecular recording; (2) \code{"nodes"} indicates that both leaves and internal nodes correspond to observations, as is typically the case for live-cell tracking; and (3) \code{"subset"} indicates that only a subset of nodes correspond to observations (see \code{pycea.datasets.packer19} for an example in which a subset of nodes in the \textit{C. elegans} lineage tree was measured). Trees are validated against the observations upon assignment, and the selected \code{alignment} option controls how trees are pruned when the object is sliced, ensuring consistency between trees and other measurements. This pruning makes slicing noticably slower than in AnnData (Figure~S1A).

A TreeData object can hold multiple trees each specified by a key (Figure~2C). \code{allow_overlap=False} allows for trees with disjoint observations, such as one tree per founding clone in a tumor. Multiple trees can also describe the same observations (\code{allow_overlap=True}), as when comparing alternative reconstructions of the same lineage data (e.g. neighbor-joining and UPGMA trees). Pycea functions operate on all trees by default or on a selected subset. Thus, multi-clone experiments and comparisons of reconstruction methods do not require separate objects.

\subsection*{Pycea provides composable tools for visualizing and analyzing lineage trees}

Pycea is organized like Scanpy, with preprocessing (\code{pp}), tools (\code{tl}), plotting (\code{pl}) and data access (\code{get}) functions that operate on TreeData objects and write their results back to \code{obs}, \code{obsp}, \code{obst}, or \code{uns}. Pycea also provides a \code{datasets} module for loading example datasets, including the mouse embryo and tumor datasets analyzed below.

Complex tree plots can be constructed through successive \code{pl} calls (Figure~2D). For example, users can first draw tree branches with \code{pycea.pl.branches}, add nodes with \code{pycea.pl.nodes}, and finally add leaf annotations with \code{pycea.pl.annotation}. Each function can color the tree by attributes such as branch length, node depth, categorical or continuous annotations, and gene expression. Trees can be drawn in rectangular or polar layouts, and the plotting functions automatically add legends for categorical variables and colorbars for continuous variables. These functions scale to large trees, rendering trees with 100,000 cells in seconds (Figure~S1B).

Pycea provides tools to quantify lineage structure and relate it to other cellular measurements (Figure~S1C). \code{tree_distance} and \code{tree_neighbors} measure lineage relationships, while \code{distance}, \code{compare_distance}, and \code{neighbor_distance} quantify and compare distances in expression or other embeddings. \code{clades} groups cells by shared ancestors, and \code{sort} orders branches by node attributes. \code{ancestral_states} infers internal node states using averaging or parsimony algorithms \citep{fitch1971,hartigan1973,sankoff1975}, while \code{parsimony} and \code{fitch_count} quantify state transitions \citep{quinn2021metastasis}. \code{ancestral_linkage} measures relatedness between cell categories, \code{autocorr} quantifies lineage autocorrelation \citep{moran1950}, and \code{partition_test} compares sibling clades. To study growth dynamics, \code{n_extant} counts lineages over time, \code{fitness} estimates relative fitness \citep{neher2014}, and \code{expansion_test} identifies clonal expansions \citep{yang2022phylodynamics}. These methods scale to trees of 100,000 cells, with runtimes of 1–20 seconds for representative operations (Figure~S1C), enabling integrated analysis of lineage, expression, and spatial data.

\subsection*{Lineage relationships and gene heritability in a mouse embryo}

To illustrate these tools, we analyzed mouse embryo lineage tracing data from Colgan et al. (\code{pycea.datasets.colgan26}) \citep{colgan2026}. We focused on a single embryo collected at E9.5 (E9.5-R1), containing 262,534 lineage-traced cells. Each cell was assigned to a germ layer and a broad lineage category, such as neural crest, neural ectoderm, or surface ectoderm (Figure~3A–C). The lineage tree contains hundreds of distinct germ layer clades, reflecting substantial proliferation before germ layer specification.

We used \code{pycea.tl.ancestral_linkage} to examine how closely related different cell populations are. For each cell, this function identifies the lowest common ancestor (LCA) shared with a target population (Figure~3D). Because this measure depends on tree topology and target population size, it is normalized by subtracting the mean under random permutation of the target cells \citep{chan2019molecular,fang2022qfm}. For example, computing the ancestral linkage of every cell to the neural crest reveals that dorsal neural ectoderm and surface ectoderm cells share recent ancestry with this population (Figure~3D). We also computed mean pairwise linkage between lineages by averaging values for each source–target pair and symmetrizing the resulting matrix (Figure~3E). This analysis groups lineages by germ layer and reveals a surprising association between the neural crest and intermediate mesoderm, potentially reflecting their similar mediolateral positions in the embryo.

Pycea also enables gene-level analysis. For example, we quantified the heritability of gene expression across the lineage tree using Moran's I statistic (\code{pycea.tl.autocorr}) \citep{moran1950}. Genes whose expression is unrelated to lineage have values near 0, whereas genes with correlated expression among closely related cells have higher values (Figure~3F). In the mouse embryo, the four most heritable genes were \textit{Hoxb9}, \textit{Meox1}, \textit{Otx2} and \textit{Sox2}, all of which are lineage-defining transcription factors or patterning genes whose expression is restricted to specific regions of the UMAP (Figure~3G).

\subsection*{Cell proliferation dynamics in a tumor metastasis model}

To demonstrate how Pycea can be used to analyze spatially resolved lineage tracing data, we analyzed a 4T1 tumor dataset from Koblan et al. (\code{pycea.datasets.koblan25}) \citep{koblan2025}. We focused on three sections from a single tumor (M3-T1), containing 145,954 cells with gene expression and lineage information measured \textit{in situ} using imaging. We used \code{pycea.tl.clades} to group cells into clades based on shared ancestry at day 5, then plotted these clades alongside the lineage marks used to reconstruct the tree (Figure~4A). Because TreeData is compatible with Squidpy, we could readily visualize the spatial distribution of these clades, revealing that spatially proximal cells are much more likely to share ancestry (Figure~4B).

We then used the tree to quantify proliferation dynamics. \code{pycea.tl.n_extant} counts the lineages alive at each time point, revealing how the relative abundance of these clades changes over time (Figure~4C). \code{pycea.tl.fitness} estimates the relative fitness of each cell from tree topology and branch lengths using selection-biased diffusion \citep{neher2014} (Figure~4D,E). Notably, the blue and purple clades tended to have higher fitness, with the highest-fitness cells occupying the tumor interior near normal lung tissue, where nutrients may be more abundant. We also investigated whether gene expression was associated with fitness by computing the Pearson correlation between each gene's expression and fitness across cells (Figure~4F). The strongest correlation was observed for \textit{Cldn4}, which is expressed in the tumor interior and associated with an epithelial-like state (Figure~4G,H) \citep{koblan2025}.
